\chapter{磁介质}
\section{磁化强度和磁化电流}
我们曾经提到磁体具有吸引铁磁性物质的能力, 并且将这种能力定义为磁性. 事实上不仅仅是磁体和铁磁性物质具有磁性, 其他物质或多或少都具有磁性. 使物质具有磁性的物理过程称为磁化, 而一切能够磁化的物质均称为磁介质. 下面我们根据安培的分子电流假说, 来介绍一些与磁介质相关的性质.

安培分子电流假说告诉我们, 已磁化物质的磁性来源于物质内部有规则排列的分子电流. 这里我们以软铁棒为例说明磁化的微观机制.
\insertfig{3-1.png}{磁化的微观机制}{0.25}

按照安培的观点, 磁棒内每个磁 ``分子''(这里泛指磁介质中的微观基本单元)都相当于一个环形电流. 在无外场的作用下, 分子电流的取向是杂乱的, 其磁矩相互抵消. 这样宏观来看铁棒就不产生磁性. 现在我们加入外加磁场$B_0$, 称之为磁化场. 在磁化场的力矩作用下, 各分子环流的磁矩在一定程度上沿着场的方向排列起来, 此时我们便说磁场被磁化了.

由图可以看出, 在均匀介质均匀磁化的时候由于分子环流的回绕方向是一致的, 在介质中任何两个分子环流中相邻的一对电流元方向总是彼此相反且抵消的, 于是只有横截面边缘的各段电流元未被抵消. 宏观看起来, 这横截面内所有分子环流的总体效果与沿着截面边缘的一个大环形电流等效.

为了描述磁介质的磁化方向和程度, 我们引
入磁化强度矢量的概念.
\begin{definition}[磁化强度矢量]
磁化强度矢量定义为单位体积内分子磁矩的矢量和, 也即:
\[
\mathbf{M} = \frac{\sum_i \mathbf{m}_i}{\Delta V}
\]
\end{definition}

也正如电介质中极化强度矢量$P$与极化电荷之间有一定关系, 磁介质中磁化强度矢量$M$与磁化电流之间也有一定的关系.

\begin{figure}[h]
  \centering
  % 第一张图片
  \begin{minipage}[t]{0.4\textwidth}
    \centering
    \includegraphics[width=\textwidth]{3-2.png}
    \caption{磁化强度与磁化电流的关系}
  \end{minipage}
  \hfill
  % 第二张图片
  \begin{minipage}[t]{0.4\textwidth}
    \centering
    \includegraphics[width=\textwidth]{3-3.png}
    \caption{左图的微观截面}
  \end{minipage}
\end{figure}

为了方便解决问题, 我们假定每个体积元内的分子磁矩都完全一致. 设单位体积内的分子环流数为$n$, 环内具有相同的电流$I$, 环具有相同的面积$a$, 于是介质中的磁化强度为
\[
M= nIa
\]
我们在磁介质中划出一个宏观的面$S$来考察其有无分子电流通过. 如图所示, 我们只需要考虑图中仅与$S$一次相交的分子环流. 首先我们取界面环线上的一个线元$\d \mathbf{l}$. 以$\d \mathbf{l}$为轴线, $\mathbf{a}$为底面做一个柱体, 其体积为$\d \mathbf{l}\cdot \mathbf{a}$. 所有中心在此柱体内的分子环流都被$\d \mathbf{l}$穿过, 每个分子环流贡献一个通过$S$面的电流$I$, 于是$\d \mathbf{l}$穿过的所有分子环流总共贡献电流为
\[ n I \mathbf{a}\cdot \d \mathbf{l}=\mathbf{M}\cdot \d\mathbf{l}\]

于是我们得到磁化强度与磁化电流$I'$的分布有如下关系:
\[
\oint_L \mathbf{M}\cdot \d \mathbf{l} = \sum_{\text{L内}} I'
\]

为了得到磁化强度与介质表面磁化电流的关系, 我们将上式应用于图示的回路. 设此介质表面单位长度上的磁化电流为$i'$, 则串过矩形回路的磁化电流为$I'=i' \Delta l$, $M$的积分只在介质表面内的一边上不为$0$, 其贡献为$M_t\Delta l$, 从而得到:$i'$, 则闯过矩形回路的磁化电流为$I'=i' \Delta l$, $M$的积分只在介质表面内的一边上不为$0$, 其贡献为$M_t\Delta l$从而得到:
\[
M_t \Delta l = l' \Delta l \Rightarrow M_t = i'
\]
考虑方向, 得到:
\[
i' = \mathbf{M} \times \mathbf{n}
\]

\insertfig{3-4.png}{磁化强度与表面磁化电流的关系}{0.25}

\section{磁介质中的磁场基本定理}
我们将产生磁化场的电流称为传导电流, 与磁化电流相区别. 如果已知磁介质的磁化强度分布, 那么我们就可以根据前一节的内容确定介质内和表面的磁化电流分布. 设由传导电流和磁化电流产生的磁感应强度分别是$\mathbf{B}_0$和$\mathbf{B'}$, 则总磁感强度为二者之和, 即:
\[
\mathbf{B} = \mathbf{B}_0+\mathbf{B'}
\]
这两种磁场都由毕奥-萨伐尔定律所决定, 磁介质的全部作用在于提供磁化电流作为新的场源. 显然$\mb{B}$应该满足真空中静磁场的高斯定理和安培环路定理:
\[
\oiint_S \mathbf{B} \cdot \d \mathbf{S} = 0
\]
\[
\oint_L \mathbf{B} \cdot \d \mathbf{l} = \mu_0\sum(I_0+I')
\]

也就是说磁介质中仍满足高斯定理. 由此, 我们在两介质表面做一个柱形的高斯面, 其上下底面分别位于两介质中. 当所取高斯面厚度趋于$0$时, 其侧面上的磁通量趋于$0$. 于是上下底面的通量积分和为$0$. 因为所取高斯面为柱面, 上下底面面积相同, 于是得到:$0$时, 其侧面上的磁通量趋于$0$. 于是上下底面的通量积分和为$0$. 因为所取高斯面为柱面, 上下底面面积相同, 于是得到:
\[
B_{1n}=B_{2n}
\]
也即: \textbf{在介质表面磁感应强度法向连续.}

仿照电介质中我们引入辅助矢量消去极化电荷, 我们同样可以引入一个辅助矢量来消去磁化电流.

\[
\oint_L \mathbf{B} \cdot \d \mathbf{l} = \mu_0\sum I_0+\mu_0\sum I'
\]
用磁化强度替换掉磁化电流:
\[
\oint_L \mathbf{B} \cdot \d \mathbf{l} = \mu_0\sum I_0+\mu_0 \oint_L M \cdot \d \mathbf{l}
\]
两边同除$\mu_0$, 移项, 得到:
\[
\oint_L \left(\frac{\mb B}{\mu_0}-\mb M\right)\cdot\d\mb l=\sum I_0
\]

于是我们引入:
\begin{definition}[磁场强度]


\[
\mathbf{H} = \frac{\mb{B}}{\mu_0}-\mb{M}
\]
称之为磁场强度.
\end{definition}
实验表明, 对于各项同性的线性磁介质, 磁化强度与磁场强度有简单正比关系:
\[
\mathbf{M} = \chi_m \mathbf{H}
\]

代入前式, 我们得到:
\[
\mathbf{B} = \mu_0(1+\chi_m)\mathbf{H}
\]
取相对磁导率$\mu_r$和磁导率$\mu$:
\[
\mu_r = 1+\chi_m,\; \mu = \mu_0 \mu_r
\]
则在各向同性磁介质中, 我们有:
\[
\mathbf{B} = \mu_0\mu_r \mathbf{H} = \mu \mathbf{H}
\]

处理磁介质问题的思路是与电介质类似的, 利用我们已经知道分布的传导电流先求解相关量, 最后给出相关的结果. 下面看例题:
\begin{example}

一半径为$R_1$的长直导线上包裹着相对磁导率为$\mu_r>1$的磁介质, 外半径为$R_2$. 导线中通有电流$I$. 求磁介质内部磁场的磁感应强度大小和磁介质界面处磁化电流的大小和方向.

首先由磁介质的安培环路定理:

\[
\oint_L \mathbf{H}\cdot \d \mb{l} = I \Rightarrow 2\pi r H = I (r>R_1) \Rightarrow H = \frac{I}{2\pi r}
\]

由线性介质关系:
\[
B = \mu H = \frac{\mu_0\mu_r I}{2\pi r}, \; M = \chi_m H =\frac{(\mu_r-1) I}{2\pi r}
\]

由面电流密度与磁化强度的关系, 内表面面磁化电流密度方向与传导电流相同:
\[
\alpha_1' = M(R_1) = \frac{(\mu_r-1) I}{2\pi R_1}
\]

内表面总的电流大小为:
\[
I_1' = 2\pi R_1\alpha_1' = (\mu_r-1) I
\]

同理, 外表面面磁化电流密度与传导电流方向相反:
\[
\alpha_2' = M(R_2) = \frac{(\mu_r-1) I}{2\pi R_2}
\]

外表面总的电流大小为:
\[
I_2' = 2\pi R_2\alpha_2' = (\mu_r-1) I
\]
\end{example}


\section{磁介质的磁化规律和机理}
对于电介质来说, 一般$\chi_e>0$. 但磁介质的情况较电介质更加复杂. 磁介质大致可以分类为顺磁质, 抗磁质和铁磁质三种. 对于顺磁质, $\chi_m>0$, 对于抗磁质,$\chi_m<0$.

先来研究磁介质分子电流的微观机制. 我们考虑电子绕核做半径为$r$的圆周运动. 电子的角动量为
\[
\mathbf L=\mathbf r\times m\mathbf v=mvr\hat{\mathbf k}
\]
单位时间内通过轨道横截面的电量为电流, 其大小为:
\[
I=\frac{ev}{2\pi r}
\]
于是因为轨道运动带来的磁矩为:
\[
\mathbf{m}_{orb} = -I\pi r^2\hat{\mathbf k}=-\frac{1}{2}evr\hat{\mathbf k}
\]
于是电子的轨道磁矩可以写为:
\[
\mathbf{m}_{orb}=-\frac{e}{2m}\mathbf{L}
\]
一般我们定义磁矩与角动量的比值为旋磁比:
\[
\gamma_{orb}=-\frac{e}{2m}
\]

以上的磁矩由电子的轨道运动带来. 电子同时具有一种叫做自旋角动的内禀性质, 与电子的运动情况没有关系, 是电子的固有属性. 并且, 自旋角动量找不到任何的经典运动对应, 也就是说自旋角动量不仅并非电子绕着自身旋转带来的角动量, 而且是一个纯粹的量子效应量. 其旋磁比为:
\[
\gamma_{spin}=-\frac{g_s e}{2m}\simeq-\frac{e}{m},\qquad \mathbf m_{spin}=\gamma_{spin}\mathbf S\quad(g_s\simeq2.0023)
\]
其中$S$是自旋角动量.

顺磁性物质的磁化机制大致可以描述如下: 在外加磁场的作用下, 原来因为热运动而完全随机取向的分子磁矩倾向于沿着磁场方向排列, 从而产生沿着某方向的磁矩. 然而由于热运动的存在, 这种排列不可能完全整齐. 外磁场越强, 温度越低, 排列越整齐, 顺磁效应也就越强. 顺磁物体分子磁矩有序排列之后产生的磁矩与外加磁场方向相同, 该磁矩产生的磁场方向业也与外加磁场方向相同, 于是宏观上看总磁场被增强.

根据上面的机制介绍, 我们可以看出温度越高, 顺磁效应越弱. 同时当我们撤去外加磁场时, 磁化会立刻消失.

抗磁性物质的磁化机制大致可以描述如下: 在有外加磁场的情况下, 电子轨道会感应出一个与外磁场反向的磁场, 从而产生抗磁性. 由此可见, 抗磁性是普遍存在的. 同时由于其由电子轨道运动引起, 抗磁性通常非常弱, 但不受温度影响.

在对抗磁物质施加外场时, 其会受到力矩, 并使电子的轨道运动角动量产生影响:
\[
\mathbf{M} = \mb{m}_{orb}\times \mathbf{B}_0 \Rightarrow \d \mathbf{L} = \mathbf{M} \d t = (\mathbf{m}_{orb}\times \mb{B}_0)\d t
\]
带入电子的磁矩与角动量之间的表达式, 我们得到:
\[
\d \mathbf{L} = -\frac{e}{2m}(\mathbf{L}\times\mathbf{B}_0) \d t
\]
从而电子将以外磁场方向为轴产生进动. 这种进动的结果类似于一个反向的环形电流, 将产生一个反向的附加磁矩.

上述进动本身只改变原有磁矩的方向. 抗磁响应来自外场对电子轨道运动的微小修正: 其诱导磁矩与外场反向. 经典拉莫尔模型可用于说明符号和量级, 但实际原子的抗磁性需要量子力学描述. 这里不把单个电子的轨道频率变化当作一般材料磁化率的定量推导.

以上两种磁介质的磁性其实都很弱, 它们的$|\chi_m|\ll 1$, 而且往往都是常数.

对于铁磁体, $\mathbf{M}$和$\mathbf{H}$的关系不再一一对应, 还与铁磁体的磁化历史有关. 除此之外, 铁磁体的磁化强度随外加磁场的变化还存在滞后性, 称为磁滞效应.

对于铁磁质, 其磁化强度往往远大于磁场强度, 因此:
\[
B = \mu_0(H+M) \approx \mu_0 M
\]
也就是说, $B\sim H$曲线与$M\sim H$曲线相似.

下面我们介绍测定铁磁物体极化规律的方法.
\insertfig{3-5.png}{铁磁质的磁滞回线}{0.3}

我们取一块铁磁质, 其从未被磁化, 从$0$开始, 逐步地增大铁磁体内的磁场强度$H$, 这各过程在图纸中表现为$OS$, 该曲线称起始磁化曲线. 该曲线显示铁磁质在一开始的磁化过程中, 磁感应强度$B$也会非线性地增大. 在磁感应强度达到饱和之后, 逐步减小磁场强度, 其磁感应强度并不会沿着$SO$减小, 而是会沿着$SR$曲线减小. 当$H$减小到$0$时, 磁感应强度并不减小到$0$. $R$点处对应的磁感应强度$B_R$称为剩余磁感应强度. 进一步加入反向的磁场强度, 当磁场强度达到$C$点对应值时, 磁感应强度减小为$0$.$C$点对应的$H_C$称为矫顽力. 继续反向增大$H$, 磁感应强度沿着$CS'$曲线达到反向饱和, 此后反向增大磁场强度, 磁感应强度沿着$SS'$变化直到再次达到饱和, 上述曲线称为磁滞回线.

铁磁质的磁性主要来源于电子的自旋磁矩, 在没有外场的作用时, 铁磁质中的电子自旋将在小范围内自发有序排列, 形成一个个小的自发磁化区, 这样的自发磁化区称为磁畴. 量子力学理论告诉我们, 这种自发磁化来源于电子之间的 ``交换作用'', 使得电子自旋平行排列时能量更低. 这种作用是纯量子效应, 没有经典理论的对应.

在未磁化的铁磁质中, 各磁畴内的自发的磁化方向不同, 从而在宏观上不显示出磁性. 在加外磁场后, 起初磁化方向与磁化场方向相近的那些磁畴将会把临近的磁化方向与磁化方向相反磁畴的领域吞并一些, 从而显示出宏观的磁性. 当所有的磁畴都沿着磁化场的方向排列好时, 介质磁化达到饱和. 因此饱和磁化强度就等于每一个磁畴中原有的磁化强度. 当所有磁畴中的元磁矩都完全排列好时, 产生的磁化强度非常大. 这就是铁磁质的磁性很强的原因. 介质中的杂质和内应力在外场去掉后阻碍磁畴恢复原有状态, 这导致了磁滞现象的产生.

